\section{Scope, source baseline, and conventions}
\label{bhr:sec:scope}

The starting point is the ProveIt repository at commit
\snapshot, whose recorded time is 11 October 2026, 01:57:13 UTC.
We obtained the canonical manuscript's 81 chapter files, its README,
validation record, preamble and references, and the text members of all
eleven archives then present in \src{docs/incoming}. The source review
was targeted at the identities used below; it was not a line-by-line
soundness audit of every volume. The accompanying source inventory
identifies the snapshot and the incoming Git blob identifiers.

The newest source, \emph{Harmonic Parity and Resolvent Identities}
\cite{bhr:Parity}, already settles the rational cotangent-resolvent
problem, the balanced negapolygamma ladder, and a quartic directional
normal form. Those conclusions are baseline results here. Its
questions R4 and R6 instead ask for nonintegral complex kernel powers
and individual multiple Gamma transforms. The cubic harmonic report
\cite{bhr:Cubic} asks for a further bridge from higher digamma powers to
harmonic Laurent coefficients. These are the principal targets of
the present continuation.

\begin{table}[htbp]
\small
\begin{tabularx}{\textwidth}{@{}>{\raggedright\arraybackslash}p{0.23\textwidth}Y@{}}
\toprule
Source target & Exact scope pursued in this article\\
\midrule
Parity R6: individual Barnes functions & A finite formula for every
polynomially weighted, unit-period multiple Gamma resolvent, including
all normalization polynomials; an explicit $G$ formula.\\[5pt]
Parity R4: complex powers & A joint spectral transform and conversion
at positive integer exponent crossings, with the branch fixed along
the integration interval.\\[5pt]
Cubic Q6: higher digamma powers & The fourth-power finite part in
ordinary zeta/Stieltjes data and explicitly defined depth-two and
depth-three harmonic Laurent coefficients; an absolutely convergent
series representation.\\[5pt]
Parity R10: a transverse cyclic coordinate & A convergent integral
representative and its use in the fifth cyclic directional layer.\\
\bottomrule
\end{tabularx}
\caption{Developments relative to the inspected reports. Theorems below
state the precise domains; a reduction to higher-depth data is not a
reduction to an unspecified algebra of familiar constants.}
\label{bhr:tab:targets}
\end{table}

\subsection{What remains a conjecture}

The short Gaussian $S_6$ evaluation with label
\src{cycloquot:conj:S6} and the revised $S_8$ evaluation with label
\src{s8new:conj:S8} remain conjectural. The older frozen $S_8$ vector
was rigorously rejected in the existing manuscript; it is a different
candidate. None of the transformations in this article is an exact
certificate for either current candidate. Similarly, the familiar
cubic constants are not declared to belong to a smaller classical
constant algebra simply because an integral or a new convergent
series represents them.

\subsection{Notation}

We use the standard Hurwitz Laurent convention
\begin{equation}\label{bhr:eq:stieltjes-convention}
 \zeta(1-s,x)=-\frac1s+
    \sum_{m=0}^{\infty}\gamma_m(x)\frac{s^m}{m!},
 \qquad \gamma_0(x)=-\psi(x),\quad \psi=\frac{\Gamma'}\Gamma.
\end{equation}
Thus $\gamma_m=\gamma_m(1)$ and $\gamma=\gamma_0$.
An ordinary prime on $\zeta$ means a derivative in its first,
spectral argument. A superscript on $\psi$ or $\gamma_m(x)$ means
an argument derivative. We write
\[
 \Li_s^{[j]}(q)=\partial_s^j\Li_s(q),\qquad
 H_n^{(r)}=\sum_{k=1}^n k^{-r},\qquad H_n=H_n^{(1)}.
\]
All logarithms of positive real numbers are real. Complex logarithms
are fixed explicitly whenever they occur in a power. Products of
positive real Gamma and Barnes functions on $(0,1)$ use their real
logarithms.

The endpoint finite part in this article is
\begin{equation}\label{bhr:eq:fp-definition}
 \FP_x\int_0^1 f(x)\,dx
   = [\delta^0(\log\delta)^0]
       \int_{\delta}^{1-\delta}f(x)\,dx.
\end{equation}
This definition applies when the integral has a finite asymptotic
expansion of nonvanishing powers and logarithms plus a constant and a
vanishing remainder. For complex exponents, the removed powers may
oscillate as well as diverge. Our proofs either establish that expansion or
give an equivalent explicit subtraction. A parameter Laurent
constant, denoted $\CT$, is a different operation. Their equality
is asserted only after the relevant correction has been proved.

\subsection{Classical inputs and the meaning of a contribution}

The Hurwitz Fourier formula, the argument-shift law, Gamma
derivatives, Bernoulli polynomial evaluations, and ordinary
polylogarithm calculus are standard \cite{bhr:DLMFZ,bhr:DLMFG}.
Espinosa--Moll already developed Hurwitz integrals and balanced
negapolygamma functions \cite{bhr:EMIntegrals,bhr:EMGeneralized}.
The multiple-Gamma/Hurwitz-zeta conversion is classical
\cite{bhr:Adamchik,bhr:Vardi}. The Fourier coefficients of
$x\log\Gamma(x)$ and $\log G(x)$ are also known, in particular in
Mez\H{o}'s paper \cite{bhr:Mezo}. We do not reclassify any of those
facts as discoveries.

Our Barnes result is a finite resolvent reduction, valid with an
arbitrary polynomial weight and every specified normalization
constant. Its spectral form supplies an all-index Stieltjes
extension at the same time. The later sections isolate the precise
additional harmonic and transverse coordinates rather than hiding
them inside unnamed finite parts. These are developments relative
to the pinned source programme. Numerical checks in the package
support independent error detection; the proofs in the text carry
the mathematical claims.
